\chapter{Carotid Endarterectomy}

\section*{What Is This Surgery?}
Carotid endarterectomy (CEA) is a definitive surgical procedure aimed at removing atherosclerotic plaque from the carotid arteries, which are vital blood vessels responsible for supplying oxygenated blood to the brain. This intervention is most commonly performed at the carotid bifurcation, the point where the common carotid artery divides into the internal and external carotid arteries, as this anatomical site is particularly prone to plaque accumulation and turbulent flow. The primary objective of CEA is to restore unimpeded cerebral blood flow and, crucially, to eliminate the source of emboli (plaque fragments or thrombus) that can dislodge and travel to the brain, causing transient ischemic attacks (TIAs) or debilitating ischemic strokes. It represents a cornerstone revascularization technique for patients at significant risk of stroke due to clinically relevant carotid artery disease.

\section*{Alternative Names}
\begin{itemize}
  \item CEA
  \item Carotid artery surgery
  \item Carotid plaque removal
  \item Carotid thromboendarterectomy
  \item Open carotid surgery
  \item Endarterectomy of the carotid artery
  \item Carotid revascularization
\end{itemize}

\section*{Relevant Surgical Anatomy}
The carotid arterial system is a critical component of the cerebral circulation, originating from the brachiocephalic artery on the right and directly from the aortic arch on the left. The common carotid artery (CCA) ascends in the neck, typically bifurcating at the level of the C3-C4 vertebral bodies, or the superior border of the thyroid cartilage, into the internal carotid artery (ICA) and the external carotid artery (ECA). The carotid bulb, a slight dilation at the origin of the ICA, is a common site for atherosclerotic plaque formation and contains baroreceptors of the carotid sinus, which regulate blood pressure. The ICA ascends without branches in the neck, entering the skull to supply the anterior and middle cerebral arteries, which are major contributors to the Circle of Willis. The ECA gives off numerous branches in the neck and face, including the superior thyroid, lingual, facial, occipital, and superficial temporal arteries, and is typically ligated or clamped during CEA to prevent back-bleeding.

Key anatomical relations in the surgical field include the sternocleidomastoid muscle, which is retracted laterally to expose the carotid sheath. Within the carotid sheath lie the common carotid artery, the internal jugular vein (lateral to the artery), and the vagus nerve (posterior and medial to the artery). Other critical nerves to identify and protect include the hypoglossal nerve (CN XII), which crosses the ICA anteriorly, the superior laryngeal nerve (a branch of the vagus), the recurrent laryngeal nerve (especially on the right), the marginal mandibular branch of the facial nerve (CN VII) superiorly, and the great auricular nerve (sensory) crossing the sternocleidomastoid. The carotid body, a chemoreceptor, is located at the carotid bifurcation and is often manipulated during dissection, potentially leading to transient blood pressure fluctuations. Lymphatic drainage of the region occurs via the deep cervical lymph nodes.

\section{Surgical Principle \& Rationale}
The fundamental surgical principle of carotid endarterectomy is the mechanical removal of the atherosclerotic plaque that has accumulated within the subintimal layer of the carotid artery wall. Atherosclerosis leads to the formation of a fibrofatty plaque, which progressively narrows the arterial lumen (stenosis) and, critically, can become unstable, serving as a nidus for thrombus formation or a source of emboli. These emboli, consisting of plaque fragments or thrombus, can dislodge and travel distally into the cerebral circulation, causing transient ischemic attacks (TIAs) or ischemic strokes. The rationale for CEA is to eliminate this embolic source and to restore adequate cerebral blood flow by widening the arterial lumen.

Technically, the procedure involves making a longitudinal incision (arteriotomy) into the affected carotid artery segment, typically extending from the common carotid artery across the bifurcation into the internal carotid artery. The atheromatous plaque, which lies between the intima and media layers of the arterial wall, is then carefully dissected and peeled away (endarterectomy) from the underlying media. This process aims to leave a smooth, wide, and unobstructed lumen. To prevent immediate re-narrowing (restenosis) and to ensure a larger, more durable lumen, the arteriotomy is almost universally closed with a patch of synthetic material (e.g., Dacron, PTFE) or autologous tissue (e.g., saphenous vein). This patch angioplasty technique is crucial for long-term patency and improved outcomes compared to primary closure.

\section{History \& Surgical Evolution}
The concept of surgically intervening on the carotid artery to prevent stroke emerged in the mid-20th century, driven by a growing understanding of atherosclerosis and its cerebrovascular consequences. The first successful carotid endarterectomy was performed by Dr. Michael DeBakey in 1953 at Methodist Hospital in Houston, Texas, marking a groundbreaking achievement in vascular surgery. This pioneering operation demonstrated the feasibility of directly addressing carotid stenosis.

Early procedures were often performed without the benefit of intraluminal shunts or patch angioplasty, leading to challenges with cerebral ischemia during clamping and higher rates of restenosis. Over the subsequent decades, significant refinements in surgical technique and perioperative management dramatically improved outcomes. The routine use of intraluminal shunts, pioneered by J. Gordon McMurrich and others, became standard practice to maintain cerebral perfusion during the arterial clamping phase, thereby minimizing intraoperative stroke risk. The widespread adoption of patch angioplasty, championed by surgeons like Dr. E. Stanley Crawford, significantly reduced the incidence of restenosis and improved long-term patency. The efficacy of CEA in preventing stroke was definitively established by landmark randomized controlled trials in the 1990s, including the North American Symptomatic Carotid Endarterectomy Trial (NASCET), the European Carotid Surgery Trial (ECST), and the Asymptomatic Carotid Atherosclerosis Study (ACAS). These trials provided robust, evidence-based guidelines for patient selection, solidifying CEA's role as a standard of care for both symptomatic and selected asymptomatic patients with high-grade carotid stenosis.

\section{Clinical Indications}
Carotid endarterectomy is indicated for patients with significant carotid artery stenosis who are at increased risk of ischemic stroke. The decision for surgery is based on the degree of stenosis, the presence of symptoms, and the patient's overall surgical risk.

\begin{itemize}
  \item Symptomatic carotid artery stenosis of 50-99\% in patients who have experienced a transient ischemic attack (TIA) or minor ischemic stroke attributable to the ipsilateral stenotic lesion.
  \item Amaurosis fugax (temporary monocular blindness) caused by ipsilateral carotid stenosis, representing a form of TIA.
  \item Asymptomatic carotid artery stenosis of 70-99\% in patients with a life expectancy of at least 5 years and low surgical risk, particularly when performed by experienced surgeons with documented low complication rates.
  \item Progressive neurological deficits or recurrent TIAs despite optimal medical therapy in patients with moderate to severe stenosis.
  \item Carotid artery stenosis identified as the likely source of an embolic stroke of undetermined source (cryptogenic stroke) after thorough diagnostic workup.
  \item Patients with carotid artery stenosis undergoing other major surgical procedures, such as coronary artery bypass grafting, where the combined risk of perioperative stroke is deemed high enough to warrant prophylactic CEA.
  \item Carotid artery stenosis with features suggestive of high-risk plaque morphology (e.g., ulcerated plaque, intraplaque hemorrhage) even at lower degrees of stenosis, though this remains a subject of ongoing research.
\end{itemize}

\section*{Clinical Positioning}
Carotid endarterectomy is primarily positioned as a first-line surgical intervention for the prevention of ischemic stroke in appropriately selected patients with symptomatic, high-grade carotid artery stenosis. In this context, it is often the preferred revascularization strategy, particularly for patients deemed to be at standard surgical risk and with favorable anatomical characteristics. For asymptomatic patients with severe stenosis, CEA also serves as a primary prophylactic measure to significantly reduce the long-term risk of future stroke.

However, CEA may be considered a secondary or alternative procedure in specific clinical scenarios. For instance, if a patient presents with anatomical challenges such as a very high carotid bifurcation, a history of prior neck radiation, or previous neck surgery leading to significant scarring, carotid artery stenting (CAS) might be considered as a primary alternative due to the increased technical difficulty and risk of open surgery. Similarly, in patients with significant medical comorbidities that elevate the risk of general anesthesia and open surgery, CAS may be favored. In rare cases where CAS is attempted but fails (e.g., inability to cross the lesion, severe dissection) or results in an incomplete revascularization, CEA could then be considered as a secondary, salvage procedure. The choice between CEA and CAS is complex and requires careful consideration of patient characteristics, plaque morphology, institutional expertise, and often involves a multidisciplinary discussion.

\section*{Surgical Technique \& Procedure}
Carotid endarterectomy is a meticulous surgical procedure typically performed under general anesthesia, though regional anesthesia with local nerve blocks and continuous neurological monitoring is an alternative in some centers. The patient is positioned supine with the head turned approximately 45 degrees to the contralateral side, exposing the neck.

The procedure involves the following key steps:
\begin{enumerate}
  \item \textbf{Anesthesia and Patient Positioning:} General endotracheal anesthesia is induced. The patient is positioned supine with the head extended and turned to the side opposite the operative carotid artery. The neck and upper chest are prepped and draped in a sterile fashion.
  \item \textbf{Incision and Exposure:} An oblique incision is made along the anterior border of the sternocleidomastoid muscle, typically extending from just below the mastoid process to the level of the cricoid cartilage or slightly lower. The platysma muscle is divided, and the sternocleidomastoid muscle is retracted laterally. Careful dissection is performed to expose the common carotid artery, its bifurcation, and the proximal internal and external carotid arteries. Critical nerves such as the vagus nerve, hypoglossal nerve, and great auricular nerve are identified and meticulously protected throughout the dissection.
  \item \textbf{Systemic Anticoagulation:} Once the carotid arteries are sufficiently exposed and ready for clamping, systemic heparin is administered intravenously to achieve an activated clotting time (ACT) typically greater than 250 seconds. This prevents thrombus formation within the clamped arterial segments and distal vessels.
  \item \textbf{Arterial Clamping and Shunting:} The common, internal, and external carotid arteries are carefully clamped in a specific sequence (usually external, then common, then internal) to temporarily stop blood flow through the operative segment. In most cases, an intraluminal shunt (e.g., Javid shunt) is inserted between the common and internal carotid arteries to maintain cerebral perfusion during the endarterectomy, thereby minimizing the risk of intraoperative stroke.
  \item \textbf{Arteriotomy and Plaque Removal:} A longitudinal incision (arteriotomy) is made along the anterior wall of the common carotid artery, extending across the bifurcation and into the internal carotid artery beyond the distal extent of the plaque. The atheromatous plaque is then meticulously dissected and removed from the underlying arterial wall, ensuring a smooth transition at the distal end of the internal carotid artery to prevent a "shelf" that could lead to dissection or embolization. The external carotid artery is also endarterectomized if significant plaque is present.
  \item \textbf{Patch Angioplasty:} After complete plaque removal and thorough irrigation of the arterial lumen, the arteriotomy is closed using a patch of synthetic material (e.g., Dacron, PTFE) or autologous vein (e.g., saphenous vein, bovine pericardium). The patch is sewn into place with fine monofilament sutures, widening the artery and reducing the risk of restenosis.
  \item \textbf{Declamping and Hemostasis:} Before final closure of the patch, the clamps are carefully removed in a specific sequence (external, then common, then internal carotid) to flush any air or debris away from the brain. Meticulous hemostasis is achieved, and the wound is irrigated with saline.
  \item \textbf{Closure:} A small suction drain may be placed in the wound to prevent hematoma formation. The platysma muscle is reapproximated, and the skin is closed in layers with sutures or staples.
\end{enumerate}

\section*{Preoperative Workup \& Preparation}
Thorough preoperative workup and preparation are paramount for optimizing patient outcomes and minimizing the inherent risks associated with carotid endarterectomy.

\begin{itemize}
  \item \textbf{Comprehensive Medical Evaluation:} A detailed medical history and physical examination are performed, focusing on cardiovascular risk factors (hypertension, diabetes, hyperlipidemia, smoking), neurological status, and overall functional capacity.
  \item \textbf{Laboratory Tests:} Routine blood tests include a complete blood count (CBC), basic metabolic panel (BMP) to assess renal function and electrolytes, liver function tests, coagulation profile (PT/INR, PTT), and type and screen for potential blood transfusion.
  \item \textbf{Cardiac Clearance:} An electrocardiogram (ECG) is standard. Patients with a significant cardiac history, symptoms of angina, or abnormal ECG findings may require further cardiac evaluation, such as an echocardiogram, stress test (pharmacological or exercise), or even coronary angiography, to ensure they can tolerate the physiological stress of surgery.
  \item \textbf{Imaging Studies:} Carotid duplex ultrasound is the primary non-invasive diagnostic tool to assess the degree of stenosis, plaque morphology, and flow characteristics. Computed tomography angiography (CTA) or magnetic resonance angiography (MRA) of the neck and brain are often performed to provide detailed anatomical information, assess plaque characteristics (e.g., ulceration, calcification), and evaluate the intracranial circulation and Circle of Willis patency.
  \item \textbf{Medication Management:} Patients are typically advised to continue antiplatelet medications (e.g., aspirin, clopidogrel) preoperatively, often with a loading dose if not already on therapy. Anticoagulants (e.g., warfarin, DOACs) usually need to be held or bridged according to specific institutional protocols. Blood pressure medications, statins, and diabetic medications should generally be continued and optimized.
  \item \textbf{NPO Status:} Patients are instructed to be NPO (nil per os) for at least 6-8 hours prior to surgery to reduce the risk of aspiration during anesthesia induction.
  \item \textbf{Prophylactic Antibiotics:} Intravenous broad-spectrum antibiotics (e.g., cefazolin) are administered within 60 minutes prior to skin incision to reduce the risk of surgical site infection.
  \item \textbf{Informed Consent:} A detailed discussion with the patient and their family regarding the procedure, its anticipated benefits, potential risks (including stroke, nerve injury, bleeding), and alternative treatment options (e.g., medical management, carotid artery stenting) is conducted, and informed consent is obtained.
\end{itemize}

\section*{Contraindications \& Precautions}
\textbf{Absolute Contraindications:}
\begin{itemize}
  \item Complete occlusion of the internal carotid artery with no distal flow on imaging, as there is no lumen to revascularize, and the risk of reperfusion injury to an ischemic brain is high without benefit.
  \item Severe, irreversible neurological deficit from a completed stroke (e.g., dense hemiplegia, global aphasia), where revascularization is unlikely to improve neurological function and may significantly increase the risk of hemorrhagic conversion of the infarct.
  \item Severe uncorrectable comorbidities that confer an unacceptably high surgical and anesthetic risk, such as severe heart failure (NYHA Class IV), unstable angina, recent myocardial infarction (within 4-6 weeks), severe uncompensated pulmonary disease, or very limited life expectancy (less than 1-2 years).
  \item Acute intracranial hemorrhage or other acute intracranial pathology that would preclude systemic anticoagulation or significantly increase surgical risk.
\end{itemize}

\textbf{Relative Contraindications and Precautions:}
\begin{itemize}
  \item Prior neck radiation therapy or previous neck surgery, which can lead to significant scarring, tissue fibrosis, and distortion of anatomical planes, making surgical dissection more challenging and increasing the risk of nerve injury or bleeding.
  \item Recurrent carotid stenosis after a previous endarterectomy, which often presents with dense scar tissue and may be managed with carotid artery stenting or repeat CEA with increased technical difficulty.
  \item Contralateral carotid artery occlusion, which places the patient at higher risk for intraoperative stroke due to reliance on the operative side for cerebral perfusion during clamping. Careful monitoring of cerebral perfusion is critical.
  \item Tandem lesions, meaning significant stenosis in both the carotid artery and the intracranial arteries, requiring careful assessment of overall cerebral blood flow and potential staged interventions.
  \item Unstable angina or recent myocardial infarction (within 3-6 months), necessitating careful cardiac risk stratification and optimization by a cardiologist prior to surgery.
  \item Severe chronic kidney disease or other systemic conditions that significantly increase anesthetic or surgical risk, requiring multidisciplinary input.
  \item Very high or very low carotid bifurcation, which can make surgical exposure and safe clamping more difficult, potentially requiring specialized techniques or alternative approaches.
  \item Significant calcification of the carotid artery, which can make arteriotomy and plaque removal more challenging and increase the risk of arterial wall injury.
\end{itemize}

\section*{Complications \& Risks}
While carotid endarterectomy is generally a safe and effective procedure, it carries potential risks and complications, some of which can be serious. These can be broadly categorized into general surgical risks and procedure-specific risks.

\begin{itemize}
  \item \textbf{General Surgical Complications:}
    \begin{itemize}
      \item \textit{Bleeding/Hematoma:} Accumulation of blood in the surgical site, which can cause swelling, pain, and potentially airway compromise, sometimes necessitating urgent re-exploration.
      \item \textit{Wound Infection:} Infection of the surgical incision, though relatively uncommon with prophylactic antibiotics.
      \item \textit{Anesthesia Complications:} Risks associated with general anesthesia, including allergic reactions, respiratory depression, cardiac events, and adverse drug reactions.
      \item \textit{Deep Vein Thrombosis (DVT) and Pulmonary Embolism (PE):} Risk of clot formation in leg veins, potentially traveling to the lungs.
    \end{itemize}
  \item \textbf{Procedure-Specific Complications:}
    \begin{itemize}
      \item \textit{Stroke:} The most feared complication, occurring in 2-3\% of cases for symptomatic patients and <1\% for asymptomatic patients in experienced centers. This can be ischemic (due to embolization of plaque/thrombus, hypoperfusion during clamping, or technical defect) or hemorrhagic (due to reperfusion injury or uncontrolled hypertension).
      \item \textit{Myocardial Infarction (MI):} Patients undergoing CEA often have underlying cardiovascular disease, making them susceptible to perioperative MI, which is a leading cause of morbidity and mortality.
      \item \textit{Cranial Nerve Injury:} Temporary or permanent damage to nerves in the neck can occur due to direct trauma, stretching, or ischemia during dissection.
        \begin{itemize}
          \item \textit{Vagus nerve (CN X):} Can cause vocal cord paralysis, leading to hoarseness, dysphagia, or aspiration.
          \item \textit{Recurrent laryngeal nerve:} A branch of the vagus, also causing hoarseness.
          \item \textit{Hypoglossal nerve (CN XII):} Can cause tongue deviation towards the injured side, affecting speech and swallowing.
          \item \textit{Marginal mandibular branch of facial nerve (CN VII):} Can cause weakness or paralysis of the lower lip, leading to an asymmetrical smile.
          \item \textit{Glossopharyngeal nerve (CN IX):} Less common, can affect swallowing and gag reflex.
          \item \textit{Great auricular nerve:} A sensory nerve, injury causes numbness in the earlobe and surrounding skin.
        \end{itemize}
      \item \textit{Hyperperfusion Syndrome:} A rare but serious complication, typically occurring in patients with chronic severe stenosis and impaired cerebral autoregulation. Characterized by severe headache, seizures, and potentially intracranial hemorrhage, usually within the first few days post-op.
      \item \textit{Restenosis:} Re-narrowing of the carotid artery at the surgical site, which can occur months or years after the procedure, sometimes requiring repeat intervention (CEA or stenting).
      \item \textit{Pseudoaneurysm:} A false aneurysm formation at the arteriotomy site, typically due to suture line failure, requiring surgical repair.
      \item \textit{Carotid Body Dysfunction:} Manipulation of the carotid body during dissection can lead to temporary or persistent blood pressure instability (hypertension or hypotension) due to altered baroreceptor function.
      \item \textit{Cerebral Edema:} Swelling of the brain, particularly after reperfusion of a chronically ischemic area.
    \end{itemize}
\end{itemize}

\section*{Postoperative Recovery Timeline}
Immediately following carotid endarterectomy, the patient is transferred to a post-anesthesia care unit (PACU) or a specialized intensive care unit (ICU) or stroke unit for close and continuous monitoring. Critical parameters include frequent neurological assessments (every 15-30 minutes initially) to detect any signs of stroke or nerve injury, continuous monitoring of vital signs, and strict blood pressure control. Maintaining blood pressure within a narrow, surgeon-specified range is paramount to prevent complications such as hyperperfusion syndrome (from hypertension) or cerebral hypoperfusion (from hypotension). Pain management is initiated, and the surgical incision is observed meticulously for signs of bleeding, hematoma formation, or swelling that could compromise the airway.

Within the first 24-48 hours, if the patient remains neurologically stable, blood pressure is well-controlled, and there are no signs of significant complications, they are typically transferred to a step-down unit or a regular surgical ward. Early ambulation is encouraged to prevent venous thromboembolism, and a soft diet is usually tolerated once the patient is fully awake and has no signs of dysphagia. The surgical drain, if placed, is generally removed within 24-48 hours. Most patients are discharged home within 1-3 days post-surgery. Upon discharge, patients receive detailed instructions on wound care, activity restrictions (avoiding heavy lifting or strenuous activities for approximately 4-6 weeks), and medication management, including continued antiplatelet therapy, statins, and optimized blood pressure control. Full recovery and a return to normal daily activities typically take 4-6 weeks, with gradual improvement in energy levels and resolution of neck discomfort.

\section*{Surgical Findings \& Pathology}
During carotid endarterectomy, the surgeon meticulously documents several key findings. These include the precise location and extent of the atherosclerotic plaque, its morphology (e.g., smooth, ulcerated, calcified, friable), the degree of luminal stenosis, and the presence of any associated thrombus. The quality of the arterial wall, the presence of any dissection planes, and the adequacy of the endarterectomy plane are also noted. The patency of the internal and external carotid arteries, as well as the back-bleeding from the internal carotid artery (an indicator of collateral cerebral flow), are assessed. Any intraoperative complications, such as nerve injury or difficulty with shunting, are also recorded.

The resected atherosclerotic plaque specimen is sent for histopathological examination. The pathology report typically confirms the diagnosis of atherosclerosis, describing the plaque's composition, which may include lipid deposits, fibrous tissue, calcification, intraplaque hemorrhage, and cholesterol clefts. The presence of ulceration on the plaque surface, which is a significant risk factor for embolization, is specifically noted. The report also confirms the completeness of the endarterectomy by assessing the margins of the resected tissue. These findings provide crucial information regarding the nature of the disease and help correlate with the patient's clinical presentation and future risk stratification.

\section{Expected Postoperative Course}
A normal and successful postoperative course following carotid endarterectomy is characterized by a smooth recovery without significant complications. Patients can expect some degree of neck pain and discomfort at the incision site, which is typically well-managed with oral analgesics and gradually resolves over days to weeks. Swelling and bruising around the incision are common and usually subside within the first week. Neurologically, the patient should remain intact, with no new deficits such as weakness, numbness, speech difficulties, or visual changes. Any pre-operative symptoms, such as transient ischemic attacks (TIAs) or amaurosis fugax, should have resolved. Blood pressure should be stable and well-controlled, often requiring careful medication adjustments in the immediate postoperative period. Patients are expected to gradually increase their activity levels, returning to light activities within a few days and resuming most normal activities, including driving, within 2-4 weeks, with full recovery typically by 6 weeks. The incision should heal cleanly, without signs of infection or excessive scarring.

\section{Postoperative Care \& Follow-up}
Regular and structured postoperative care and follow-up are essential after carotid endarterectomy to monitor for complications, assess long-term patency, and manage cardiovascular risk factors.

\begin{itemize}
  \item \textbf{Post-operative Visits:} Patients typically have a follow-up visit with the vascular surgeon within 1-2 weeks for wound check, suture or staple removal, and initial assessment of recovery. A subsequent visit may be scheduled around 1 month post-op.
  \item \textbf{Carotid Duplex Ultrasound Surveillance:} Surveillance carotid duplex ultrasound is crucial to monitor for restenosis or new disease progression in the contralateral carotid artery. It is usually performed at 1 month, 6 months, and then annually thereafter, or more frequently if clinically indicated.
  \item \textbf{Medication Management:} Ongoing management of cardiovascular risk factors is critical. This includes lifelong antiplatelet therapy (e.g., aspirin, sometimes clopidogrel), statin therapy for lipid management, and strict control of hypertension and diabetes. These medications are reviewed and adjusted as needed.
  \item \textbf{Rehabilitation:} If the patient experienced pre-existing neurological deficits or has new, mild deficits, physical, occupational, or speech therapy may be recommended to optimize functional recovery.
  \item \textbf{Activity Resumption:} Patients receive specific instructions on gradual return to normal activities, with emphasis on avoiding strenuous activities, heavy lifting, or sudden neck movements for several weeks to allow for wound healing and prevent complications.
  \item \textbf{Warning Signs:} Patients are thoroughly educated on warning signs of potential complications, such as new neurological symptoms (e.g., sudden weakness, numbness, vision changes, speech difficulties), severe or worsening headache, significant neck swelling, redness, discharge from the wound, or fever. They are instructed to seek immediate medical attention if any of these occur.
\end{itemize}

\section*{Epidemiology}
Carotid artery stenosis is a prevalent manifestation of systemic atherosclerosis, primarily affecting older adults. Its incidence and prevalence increase significantly with age, affecting approximately 5-10\% of individuals over 65 years old and up to 20\% of those over 80. Men generally exhibit a higher incidence of carotid stenosis than women, particularly at younger ages, although the gap narrows in older age groups. Carotid endarterectomy remains one of the most frequently performed vascular surgical procedures globally. In the United States, approximately 100,000 CEAs are performed annually, though this number has fluctuated over time due to evolving guidelines, improved medical therapies, and the increasing use of carotid artery stenting (CAS). The primary indication for CEA is symptomatic stenosis, accounting for a significant proportion of procedures, followed by selected cases of asymptomatic severe stenosis. Perioperative mortality and stroke rates for CEA have steadily declined over the decades due to advancements in surgical techniques, anesthetic management, and patient selection. Current rates in experienced, high-volume centers are typically below 3\% for symptomatic patients and below 1\% for asymptomatic patients, reflecting the procedure's safety and efficacy when performed by skilled teams.

\section*{Outcomes \& Prognosis}
The outcomes of carotid endarterectomy are generally excellent, particularly when performed in high-volume centers by experienced surgeons with low complication rates. For symptomatic patients with high-grade (70-99\%) carotid stenosis, landmark trials such as NASCET (North American Symptomatic Carotid Endarterectomy Trial) and ECST (European Carotid Surgery Trial) definitively demonstrated that CEA, when combined with optimal medical therapy, significantly reduces the 2-year risk of ipsilateral stroke by approximately 50-65\% compared to medical management alone. For selected asymptomatic patients with severe (70-99\%) stenosis, the ACAS (Asymptomatic Carotid Atherosclerosis Study) and ACST (Asymptomatic Carotid Surgery Trial) showed that CEA reduces the 5-year risk of ipsilateral stroke by about 50\% compared to medical therapy alone.

Functional outcomes are typically very good, with most patients returning to their baseline neurological status or experiencing resolution of pre-operative symptoms. Long-term survival after CEA is primarily influenced by the patient's underlying cardiovascular comorbidities (e.g., coronary artery disease, peripheral artery disease) rather than the CEA itself, as these patients often have diffuse atherosclerotic disease. Restenosis, defined as a significant re-narrowing of the carotid artery at the surgical site, occurs in 5-15\% of patients within 5 years, with symptomatic restenosis being less common and often amenable to repeat intervention. Prognostic factors influencing outcomes include the patient's age, severity of comorbidities (especially cardiac disease and diabetes), plaque morphology (e.g., ulceration, intraplaque hemorrhage), and the experience and volume of the surgical team. Overall, CEA remains a highly effective and durable intervention for stroke prevention in appropriately selected patients, offering a significant long-term benefit.

\section*{References}
\begin{enumerate}
  \item MedlinePlus. Carotid Endarterectomy. U.S. National Library of Medicine; 2024. Available from: \url{https://medlineplus.gov/ency/article/002940.htm}
  \item Sabiston Textbook of Surgery: The Biological Basis of Modern Surgical Practice. 22nd ed. Elsevier; 2024.
  \item American Heart Association/American Stroke Association. Guidelines for the Management of Patients with Carotid Artery Disease. Stroke. 2024;55(1):e1-e100.
  \item Rutherford's Vascular Surgery and Endovascular Therapy. 10th ed. Elsevier; 2024.
  \item North American Symptomatic Carotid Endarterectomy Trial Collaborators. Beneficial effect of carotid endarterectomy in symptomatic patients with high-grade carotid stenosis. N Engl J Med. 1991;325(7):445-453.
\end{enumerate}

\clearpage